% ==========================================
% BẢNG PHÂN CÔNG CÔNG VIỆC
% ==========================================
\section*{BẢNG PHÂN CÔNG CÔNG VIỆC VÀ MỨC ĐỘ ĐÓNG GÓP}
\addcontentsline{toc}{section}{Bảng phân công công việc}

\renewcommand{\arraystretch}{1.5}
% Thay tabular bằng longtable
\begin{longtable}{|p{0.05\textwidth}|p{0.18\textwidth}|>{\linespread{1.3}\selectfont\arraybackslash}p{0.6\textwidth}|p{0.08\textwidth}|}
    
    % --- 1. TIÊU ĐỀ Ở TRANG ĐẦU TIÊN ---
    \hline
    \centering \textbf{STT} & \textbf{Họ tên / MSSV} & \textbf{Nhiệm vụ đảm nhiệm chi tiết} & \textbf{Đóng góp} \\ \hline
    \endfirsthead
    
    % --- 2. TIÊU ĐỀ LẶP LẠI Ở CÁC TRANG SAU ---
    \hline
    \centering \textbf{STT} & \textbf{Họ tên / MSSV} & \textbf{Nhiệm vụ đảm nhiệm chi tiết} & \textbf{Đóng góp} \\ \hline
    \endhead
    
    % --- 3. CHÂN BẢNG Ở CÁC TRANG BỊ NGẮT ---
    \hline
    \endfoot
    
    % --- 4. CHÂN BẢNG Ở TRANG CUỐI CÙNG ---
    \hline
    \endlastfoot

    % ================= NỘI DUNG BẢNG =================
    
    % Thành viên 1
    \centering 1 & 
    24120018 \newline \textbf{Đào Thị Quỳnh Anh} & 
    \vspace{-0.4cm}
    \begin{itemize}[itemsep=0.15cm, topsep=0cm, partopsep=0cm]
        \item Biên soạn \texttt{part1\_demo.ipynb} minh họa trực quan kết quả Phần 1.
        \item Xây dựng \texttt{data\_gen.py}: Viết script sinh ma trận ngẫu nhiên SPD (well-conditioned) và ma trận Hilbert $H_n$ (ill-conditioned).
        \item Phân tích tính ổn định: Tính số điều kiện $\kappa_2(A)$, thu thập sai số tương đối và lý luận giải thích hiện tượng theo Định lý 3.1.
    \end{itemize} \vspace{-0.2cm} & 
    \centering 100\% \tabularnewline \hline
    
    % Thành viên 2
    \centering 2 & 
    24120347 \newline \textbf{Phan Lê Đăng Khoa} & 
    \vspace{-0.4cm}
    \begin{itemize}[itemsep=0.15cm, topsep=0cm, partopsep=0cm]
        \item Cài đặt thuật toán phân rã SVD (Phần 2).
        \item Xây dựng \texttt{solvers.py}: Tích hợp phương pháp Gauss, SVD và cài đặt thuật toán lặp Gauss-Seidel (kèm hàm kiểm tra ma trận chéo trội).
        \item Xây dựng \texttt{benchmark.py}: Lên kịch bản đo thời gian thực thi (trung bình 5 vòng lặp), tính sai số tương đối và xuất dữ liệu ra định dạng JSON chuẩn.
    \end{itemize} \vspace{-0.2cm} & 
    \centering 100\% \tabularnewline \hline
    
    % Thành viên 3
    \centering 3 & 
    24120348 \newline \textbf{Hà Đăng Khôi} & 
    \vspace{-0.4cm}
    \begin{itemize}[itemsep=0.15cm, topsep=0cm, partopsep=0cm]
        \item Trực quan hóa Toán học (Phần 2): Viết kịch bản và lập trình tạo Video demo hoạt ảnh quá trình phân rã ma trận bằng thư viện Manim.
        \item Phối hợp viết báo cáo tổng kết, biên tập tài liệu và quay dựng thành phẩm nộp bài cuối cùng.
    \end{itemize} \vspace{-0.2cm} & 
    \centering 100\% \tabularnewline \hline

    % Thành viên 4
    \centering 4 & 
    24120352 \newline \textbf{Lương Trung Kiên} & 
    \vspace{-0.4cm}
    \begin{itemize}[itemsep=0.15cm, topsep=0cm, partopsep=0cm]
        \item Phối hợp cài đặt thuật toán Phép khử Gauss và ứng dụng (Phần 1).
        \item Xây dựng \texttt{analysis.ipynb} (Phần 3): Đọc file dữ liệu JSON, vẽ đồ thị log-log so sánh thời gian thực tế với đường lý thuyết $O(n^3)$.
        \item Trình bày Markdown lý thuyết Toán học (số điều kiện, Gauss-Seidel) và viết kết luận đánh giá tương quan giữa tính ổn định và chi phí tính toán.
    \end{itemize} \vspace{-0.2cm} & 
    \centering 100\% \tabularnewline \hline
    
    % Thành viên 5
    \centering 5 & 
    24120418 \newline \textbf{Nguyễn Minh Quân} & 
    \vspace{-0.4cm}
    \begin{itemize}[itemsep=0.15cm, topsep=0cm, partopsep=0cm]
        \item Phối hợp cài đặt thuật toán Phép khử Gauss và ứng dụng (Phần 1).
        \item Biên soạn Báo cáo bằng \LaTeX (\texttt{report.tex}): Dựng khung tài liệu, xử lý các môi trường công thức Toán học phức tạp cho Phần 1 và Phần 2.
        \item Trích xuất dữ liệu, chèn biểu đồ từ Notebook vào báo cáo, viết phần Mở đầu, Kết luận và rà soát toàn bộ source code trước khi nộp.
    \end{itemize} \vspace{-0.2cm} & 
    \centering 100\% \tabularnewline \hline
    
\end{longtable}